% Kahn 算法：菱形依赖 1->2, 1->3, 2->4, 3->4；初始入度 (0,1,1,2)；按 FIFO 队列得输出 1,2,3,4。
\begin{tikzpicture}[font=\figurefont\fontsize{8}{10}\selectfont,>=Stealth,
  vtx/.style={circle,draw=FigureStroke,line width=.8pt,minimum size=6.5mm,inner sep=0pt,text=FigureInk,fill=white},
  arc/.style={draw=FigureStroke,line width=.9pt,->},
  cellz/.style={fill=FigureGreenFill,inner sep=1.5pt,rounded corners=1pt},
  note/.style={font=\FigureNoteFont\fontsize{7.5}{9}\selectfont,text=FigureMuted}]
% (a)
\begin{scope}
  \node[vtx] (v1) at (0,0) {1};
  \node[vtx] (v2) at (1.2,0.8) {2};
  \node[vtx] (v3) at (1.2,-0.8) {3};
  \node[vtx] (v4) at (2.4,0) {4};
  \draw[arc] (v1)--(v2); \draw[arc] (v1)--(v3); \draw[arc] (v2)--(v4); \draw[arc] (v3)--(v4);
  \node[note,above=1pt of v1] {入度 0};
  \node[note,above=1pt of v2] {入度 1};
  \node[note,below=1pt of v3] {入度 1};
  \node[note,above=1pt of v4] {入度 2};
  \node at (1.2,-1.9) {(a) 菱形依赖图};
\end{scope}
% (b)
\begin{scope}[xshift=3.6cm,yshift=1.1cm]
  \node[anchor=west] at (0,0) {输出};
  \node at (1.45,0) {$1$}; \node at (2.0,0) {$2$}; \node at (2.55,0) {$3$}; \node at (3.1,0) {$4$};
  \node[anchor=west] at (3.6,0) {零入度队列};
  \node[note] at (2.27,0.38) {各顶点的当前入度};
  \draw[FigureStroke,line width=.5pt] (0,-0.25)--(5.2,-0.25);
  \foreach \r/\o/\a/\b/\c/\d/\q in {1/{初始}/0/1/1/2/{[1]},2/{1}/{--}/0/0/2/{[2, 3]},3/{2}/{--}/{--}/0/1/{[3]},4/{3}/{--}/{--}/{--}/0/{[4]},5/{4}/{--}/{--}/{--}/{--}/{[\,]}}{
    \pgfmathsetmacro{\y}{-0.5*\r}
    \node[anchor=west] at (0,\y) {\o};
    \node at (1.45,\y) {\a}; \node at (2.0,\y) {\b}; \node at (2.55,\y) {\c}; \node at (3.1,\y) {\d};
    \node[anchor=west] at (3.6,\y) {\q};
  }
  % 刚降为 0 的入度
  \node[cellz] at (2.0,-1.0) {0}; \node[cellz] at (2.55,-1.0) {0};
  \node[cellz] at (3.1,-2.0) {0};
  \node at (2.6,-3.0) {(b) 每次输出后的入度与队列};
\end{scope}
\end{tikzpicture}
